\section{Introduction}
\label{sec:intro}

In the space of two years, a new empirical literature has formed around a single
question: how much of open-source software is now written by AI? The answers arrive as
percentages. Some fraction of commits, of pull requests, of comments, of lines of code
is reported as machine-generated, and those fractions are then used to argue about
maintainer burden, review overhead, contribution quality and the long-run
sustainability of open-source projects. The numbers circulate quickly, because the
question matters and because the answers are alarming.

Almost none of these numbers rest on a measurement instrument that has been validated
in the domain where it is being applied.

The instruments come in three families. Zero-shot statistical detectors such as
DetectGPT \citep{mitchell2023detectgpt}, Fast-DetectGPT \citep{bao2024fastdetectgpt} and
Binoculars \citep{hans2024binoculars} score text by its likelihood under a language
model. Heuristic filters match phrase lists, formatting signatures or punctuation
patterns. Large language models are prompted to classify authorship directly. Each
family was developed and benchmarked on a particular kind of text: news articles,
student essays, Wikipedia paragraphs, and code drawn from curated programming
benchmarks. Binoculars, to take the strongest published claim, reports detection of
over \citednum{90\%} of ChatGPT-generated text at a false positive rate of \citednum{0.01\%}
\citep{hans2024binoculars}. That claim was established on essay-length prose.

The artifacts of open-source software development are not essay-length prose. In the
AIDev dataset of agent-authored contributions \citep{li2025aidev}, the median commit
message is \AidevCommitMsgMedian{} characters long and the ninetieth percentile is
\AidevCommitMsgPNinety{}. The median pull request body is \AidevBodyMedian{} characters.
These texts are short, templated, written to a genre convention, and produced by a
globally distributed contributor base writing in a second language. Every one of these
properties is a known stressor for the detectors in question. Detector accuracy degrades
sharply on short inputs, and \citet{liang2023biased} showed that GPT detectors
systematically misclassify text by non-native English writers as machine-generated.
Open-source repositories additionally carry more than a decade of pre-existing
machine-generated text --- dependency bots, release automation, templated issue
reports --- that predates large language models entirely and that no published
prevalence study has been shown to exclude reliably.

If the instruments are miscalibrated in this domain, the consequences are not evenly
distributed. A detector with an inflated false positive rate does not merely add noise:
because non-AI artifacts vastly outnumber AI ones in most historical samples, even a
small false positive rate produces a large absolute number of false detections, and
prevalence is overestimated systematically rather than randomly. The same asymmetry
means that the bias documented by \citet{liang2023biased} would fall hardest on
contributors already least well served by the ecosystem.

\subsection{Auditing the instrument rather than the phenomenon}
\label{sec:intro:inversion}

This paper inverts the question. Rather than adding another prevalence estimate, it
asks what the existing measurement instruments actually measure when applied to
open-source artifacts, and what follows for the estimates already in the literature.

The inversion is made tractable by two sources of ground truth that require no human
annotation --- an essential property, since human annotation of authorship is precisely
what is unavailable at scale, and inter-rater reliability on ``does this look
AI-written'' is exactly the thing under dispute.

\paragraph{Positives.} Agentic coding platforms leave attribution traces. When Devin,
OpenAI Codex, GitHub Copilot, Cursor or Claude Code opens a pull request under its own
account, authorship is declared by the platform rather than inferred. The AIDev
dataset \citep{li2025aidev} collects \AidevPRs{} such pull requests across
\AidevRepos{} repositories. A second declared channel, the \texttt{Co-authored-by}
trailer, appears on \AidevTrailerRate{} of the commits in that dataset.

\paragraph{Negatives.} Text written before large language models were publicly
available cannot have been produced by them. Artifacts committed to the same
repositories before the release of ChatGPT on \LlmCutoff{} therefore form a negative set
whose purity follows from chronology rather than from judgement. This yields a
counterfactual test with an unusual property: every artifact a detector flags in this
set is a false positive \emph{by construction}, with no annotation, adjudication or
appeal. The resulting false positive rate is not an estimate of detector error under
assumptions --- it is a direct measurement.

Given measured sensitivity and specificity, published prevalence figures can then be
corrected. The estimator is not new; it is standard in epidemiology, where the problem
of inferring true disease prevalence from an imperfect screening test was solved by
\citet{rogan1978estimating} nearly fifty years ago. Its application here is
straightforward, and its absence from the software engineering literature on AI
prevalence is, we argue, the central methodological gap.

\subsection{Why this matters now}
\label{sec:intro:now}

The scale of the downstream literature makes the gap consequential. The 2026 Mining
Software Repositories Mining Challenge was built on AIDev and produced \citednum{62} accepted
papers, studying merge and rejection rates, test coverage, continuous integration
outcomes, security defects, refactoring behaviour, code clones and reviewer sentiment
in agent-authored contributions. All of them take the attribution labels as given ---
reasonably so, since AIDev's labels are platform-declared. But an entire research
programme has now been built on top of authorship attribution in under two years,
while the question of what attribution-by-detection actually measures, in the cases
where declaration is absent, has gone unasked.

That gap is beginning to be noticed from one side. \citet{khosravani2026census} audit
the \emph{declaration} channels --- configuration files, commit-message markers,
author identity, bot signatures --- across \citednum{180} million repositories, and find that
single-signal detection undercounts badly: bot-account lookup alone recovers only \citednum{3.3\%}
of the Claude Code commits their multi-method approach identifies. Their finding and
ours are complementary and point in opposite directions. They show that
\emph{declaration-based} measurement misses AI contributions that are real. We ask
whether \emph{content-based} measurement invents AI contributions that are not. Between
the two, a prevalence estimate can be bounded from both sides rather than reported as a
point.

\subsection{Contributions}
\label{sec:intro:contributions}

\begin{enumerate}
  \item \textbf{A criterion-validity evaluation} of the detection approaches used in the
    open-source prevalence literature, against attribution-based ground truth, reported
    per agent and per artifact genre rather than pooled (Section~\ref{sec:validity}).
    Pooling matters more than it may appear: \AidevTopAgentShare{} of AIDev is a single
    agent, so a pooled metric is that agent's metric wearing a general label.
  \item \textbf{A counterfactual false positive measurement} on pre-ChatGPT artifacts,
    where every detection is an error by construction, stratified by genre, language,
    text length and project characteristics, and separated from the confounding
    contribution of pre-LLM automation (Section~\ref{sec:validity}).
  \item \textbf{Bias-corrected prevalence estimates} that propagate the measured error
    rates through the Rogan--Gladen estimator, reported as intervals under best- and
    worst-case detector assumptions and contrasted with the point estimates in the
    literature (Section~\ref{sec:correction}).
  \item \textbf{An interoperable provenance model} that represents authorship claims
    together with their evidence type --- self-declared trailer, platform metadata,
    agent account, or detector inference with a score --- so that an attribution can be
    inspected and queried rather than taken on faith
    (Section~\ref{sec:provenance}).
  \item \textbf{A replication package} containing the corpus specification, the detector
    harness, the provenance vocabulary and all analysis code.
\end{enumerate}

The argument of this paper is not that AI-generated content in open source is rare, nor
that the concern motivating the prevalence literature is misplaced. It is that the
field currently cannot tell the difference between a detector finding AI text and a
detector finding short text, templated text, or text written by someone whose first
language is not English --- and that this is a measurable, correctable state of affairs
rather than an inherent limitation.
